\documentclass[10pt,a4paper]{article}
\usepackage[margin=0.5in]{geometry}  % 减小页边距
% 在文档开始处添加
\usepackage{enumitem}
\setlist[itemize]{leftmargin=15pt}
% 调整行间距
\renewcommand{\baselinestretch}{0.95}  % 减小整体行间距
\setlength{\parskip}{0.3em}  % 减小段落间距
% 调整节标题前后的间距
\usepackage{hyperref}
\usepackage{enumitem}
\usepackage{xeCJK} % 支持中文
\usepackage{geometry}

% 调整页面边距
\geometry{
  left=0.8cm,   % 左边距
  right=0.8cm,  % 右边距
  top=0.7cm,    % 上边距
  bottom=0.7cm  % 下边距
}


% Hyperlink setup for blue underlined links
\hypersetup{
    colorlinks = true, % Enable colored links
    linkcolor = blue,  % Color for internal links
    anchorcolor = blue, % Color for anchor text
    citecolor = blue, % Color for citations
    filecolor = blue, % Color for file links
    menucolor = blue, % Color for menu links
    runcolor = blue, % Color for run links
    urlcolor = blue, % Color for URLs (hyperlinks)
}
\begin{document}
\pagestyle{empty} % 不显示页码

\begin{center}
{\Large \textbf{张丰毅 Maple ZHANG}}\\
\vspace{0.5em}

\href{tel:14479022370}{+1 447-902-2370} /\href{tel:8617852621070}{+86 178-5262-1070} \textbar\ \href{mailto:me.fzhang@gmail.com}{me.fzhang@gmail.com} \textbar\
\href{https://linkedin.com/in/fengyizhang-profile}{LinkedIn}  \textbar\
\href{https://onemapl3.github.io}{Github Page}
% \href{https://onemapl3.github.io/images/resume.pdf}{English Version} 
\end{center}

\vspace{0em}

\noindent\textbf{教育背景}\\[-0.2em]
\rule{\textwidth}{1pt}
\textbf{伊利诺伊大学厄巴纳-香槟分校(UIUC)} \hfill 2022.08 - 2024.12 \\
\textit{Master of Computer Science}\\
研究方向：数据挖掘、机器学习、人机交互、云网络、物联网\\
\textbf{山东大学} \hfill 2017.09 - 2021.06\\
计算机科学与技术 工学学士  总GPA: 3.82/4.0\\
相关课程: 数据结构与算法、计算机网络、数据库系统、操作系统

% \vspace{0.2em}
% \noindent\textbf{专业技能}\\[-0.2em]
% \rule{\textwidth}{0.2pt}
% \textbf{编程语言}:\textit{Python, Go,C/C++/C\#,PHP,Java,JavaScript/TypeScript,SQL/MQL,Cipher,Thrift,Solidity}\\
% \textbf{框架/SDK}:\textit{FastAPI,Langchain,MCP,Gin,GORM,gRPC,React/Native,Express,.Net,CodeIgniter,Pytorch}\\
% \textbf{软件/平台}\textit{Git,Node,MySQL,PostgreSQL,MongoDB,Matlab,Apache,Nginx,Pulsar,Redis,Consul,Postman\\Swagger,Jmeter,Jenkins,Azure,AWS,GCP,Firebase,Docker,Wireshark,Linux
% \vspace{0.2em}}
\noindent\textbf{专业技能}\\[-0.2em]
  \rule{\textwidth}{0.2pt}
  \textbf{编程语言}：\textit{Go, Python, TypeScript/JavaScript, C/C++/C\#, Java,
   PHP, SQL, Solidity, Shell}\\
  \textbf{框架/SDK}：\textit{React, Next.js, Node.js, FastAPI, Gin, React Native, Express,
  Prisma, GORM, gRPC, REST, GraphQL, PyTorch, .Net, Tailwind CSS}\\
  \textbf{云与平台}：\textit{Git, GitHub Actions, Docker, Kubernetes, AWS(ECS/Lambda/RDS), GCP, Azure, Vercel,
  PostgreSQL, MySQL, MongoDB, SQLite, Qdrant, Redis, Celery, Nginx, Pulsar, Firebase, Jenkins}\\
  \textbf{数据与质量}：\textit{Pandas, NumPy, ECharts, Cesium, H3, Jest, Playwright, Testing, Observability}\\
  \textbf{AI/Agent 工程}：\textit{Codex, Claude Code, Cursor, GitHub Copilot, OpenClaw,
  MCP, Agent Skills, Multi-Agent Orchestration, RAG, Vector Databases, Prompt Engineering, LangChain, Anthropic API, OpenAI SDK}
  \vspace{0.2em}

\noindent\textbf{工作经历}\\[-0.2em]
\rule{\textwidth}{1pt}
\textbf{天空智算 (Space AIC)} \hfill 2026.03 - 至今\\
软件工程师, 上海, 中国\vspace{-8pt}
\begin{itemize}[noitemsep,leftmargin = 15pt]
    \item 基于 React 19、TypeScript、Vite、Cesium、Zustand 与 ECharts 开发全球首个支持万星级的卫星数字孪生系统，将实验规模从 1,800 颗卫星提升至 15,000+ 颗，并保持三维轨道/姿态回放交互
    \item 将 React/Cesium 控制面接入由拓扑生成、加权路径计算与分片卫星模拟器组成的分布式 RPC 仿真后端，统一结果帧并支撑星座设计、姿态控制、干扰、热控、电源/能源、空间环境与 ns-3 实验
    \item 实现半实物仿真流程与 microVM 在线终端，通过 HTTP/WebSocket/gRPC Terminal Gateway 完成实体/路径终端路由、资源监控、多窗口管理及会话状态控制
    \item 完成 1,584 星实验的实时 Cell/Beam 波束可视化：接入 timeline/entity-timeline，展示单星 16 波束、H3 小区网格与覆盖锥体，并解决播放中覆盖滞后、实体 ID 冲突和多星凝视跳动
    \item 开发卫星场景 AI 助手：支持 SAT Agent 卫星信息问答与仿真辅助，面向数字孪生场景提供方案分析，并探索在卫星上部署边缘算力运行大模型的天地一体化架构
    \item 打通 SDT 教育版 Notebook--Python 仿真引擎--Express API--Cesium 数据链路，以 JupyterHub/Docker 提供学生隔离环境，形成“建模--编程--仿真--可视化--提交--评价”闭环；并交付 ITU Agent 的流式 SatAgent/RAG、轨道申报、分级权限与管理后台
\end{itemize}

\textbf{Nethermind} \hfill 2023.08 - 2023.11\\
区块链研发实习生, 远程\vspace{-8pt}
\begin{itemize}[noitemsep,leftmargin = 15pt]
    \item 为StarkNet区块链浏览器Voyager实现并迁移新的架构设计, 从AWS Lambda和Starknet网关迁移到AWS ECS和RPC端点, 重写2个API并绘制2个架构图
    \item 完全重构搜索模块, 根据用户输入提供适当的搜索选项, 使搜索延迟减少约1000毫秒
    \item 通过添加索引和重写PostgreSQL查询, 优化和提高了从AWS RDS的检索效率
    \item 设计了一种提示工程方法, 用于自动检测和解码Cairo类型, 以人性化格式显示交易数据, 相较传统正则匹配提高了准确性
\end{itemize}

\noindent\textbf{微软} \hfill 2022.01 - 2022.08\\
软件工程师, 苏州, 中国\vspace{-8pt}
\begin{itemize}[noitemsep,leftmargin = 15pt]
    \item 通过QR/Accessibility代码与8个RESTful API和.Net框架设计并实现无缝的PC-Phone配对流程, 将Windows OOBE期间的完成率从28.3\%提高到35.4\%
    \item 采用微服务架构、异步编程、缓存、WebSocket、长轮询和API限流, 减少服务器和数据库压力, 在720 RPS下实现97\%的API成功率
    \item 利用JWT加密证书、设备和帐户信息, 在PC和手机之间建立基于帐户的加密信任关系, 确保数据传输和验证的安全性
    \item 添加超过2,000行单元/功能/端到端/负载测试和指标遥测, 用于异常检测、分析和故障排除, 提高服务可靠性和可用性
\end{itemize}

\newpage
\noindent\textbf{滴滴出行} \hfill 2021.01 - 2021.12\\
软件开发工程师, 北京, 中国\vspace{-8pt}
\begin{itemize}[noitemsep]
    \item 将每天服务1500万用户的兴趣点(POI)服务的6个API从PHP重构为Go, 在QPS为2000时, CPU空闲时间提高到73\%(之前为36\%), TP99延迟降低51\%
    \item 使用内部工作流引擎将业务流程划分为可插拔模块, 并将8个RESTful API替换为具有服务发现(DiSF)组件的滴滴RPC(DiRPC)协议, 以更好地进行SOA治理
    \item 应用A/B测试实现金丝雀部署、配置同步和策略实验, 并采用指标组件跟踪服务上下文中生成的日志以进行故障排除
    \item 在待命期间管理22个API和16个下游服务的在线指标警报(错误率、延迟等)
\end{itemize}

\noindent\textbf{阿里云} \hfill 2020.07 - 2020.09\\
软件工程师实习生, 杭州, 中国\vspace{-8pt}
\begin{itemize}[noitemsep]
    \item 主持基于旁路嗅探的路由回溯项目, 用于SLA监控和网络故障定位
    \item 使用Scapy构造TTL字段递增的虚假源VIP数据包序列到特定端点
    \item 采用流量镜像服务进行数据包捕获和带有回调过滤功能的异步I/O
    \item 在高级抗DDoS集群上交付, 同时可通过配置文件进行弹性扩展
\end{itemize}

\vspace{1em}

\noindent\textbf{个人项目}\\[-0.2em]
\rule{\textwidth}{1pt}
\textbf{时光机 (Time Machine)} \hfill 2026.01 - 2026.03\\
基于 SecondMe 数字分身、时间记忆蒸馏/RAG 与 SSE 的 AI 情感陪伴应用\vspace{-8pt}
\begin{itemize}[noitemsep]
    \item 围绕 SecondMe 数字分身 Agent，基于 \textbf{Next.js 16}、App Router、Prisma/SQLite、OAuth2 独立完成全栈产品并部署至 Fly.io，让用户可围绕指定日期私密地重温关系记忆
    \item 对文字/图片记忆进行时间化蒸馏，提取带日期的记忆锚点并按所选时点执行 RAG，仅注入当时可知信息，避免“现在的信息”污染对过去人物的模拟
    \item 实现双 Agent \textbf{SSE} 流式架构：SecondMe 负责主对话，Anthropic Act 疗愈 Agent 同步输出结构化情绪分析与共情回应
    \item 使用 Canvas/requestAnimationFrame 构建 TimePortal 状态机，串联引导、关系选择、记忆锚点、时间穿越与实时对话，形成完整情感叙事体验
\end{itemize}

\noindent\textbf{Type Machine (Media Transcriber)} \hfill 2026\\
AI Native 多模态采集与知识处理 PWA\vspace{-8pt}
\begin{itemize}[noitemsep]
    \item 基于 \textbf{Next.js App Router} 构建移动优先 PWA，并接入 Web Share Target，使用户可从系统分享菜单直接提交截图/视频，桌面端亦支持图片、音频、视频与 URL
    \item 设计异步多模态任务流水线，编排 OCR、语音转写与多模态模型处理，将异构媒体归一为统一的可编辑笔记结构，并避免大文件处理阻塞上传
    \item 实现连续截图有序重建、重叠段落去重、界面噪音过滤、语言识别与原文保留，解决长文章逐张 OCR 后仍需手工拼接的问题
    \item 增加摘要、标签与双语输出等 AI 后处理，并持久化任务历史及搜索/编辑/复核流程，通过 Flomo API 投递完整结果，使知识可追溯、可复用
\end{itemize}

\noindent\textbf{一页诗 (OnePoem)} \hfill 2026.04\\
跨 Web 与 iOS 的诗歌排版与导出工具\vspace{-8pt}
\begin{itemize}[noitemsep]
    \item 使用 \textbf{React} 开发移动优先的诗歌编辑器，支持模板、字体、配色、版式比例、诗集与桌面小组件，将纯文本转化为可收藏、可分享的视觉页面
    \item 实现 html2canvas PNG 导出及 SVG 降级方案，保证浏览器与受限移动运行时均可生成作品图片
    \item 通过 \textbf{SwiftUI}/\textbf{WKWebView} 封装 iOS 应用，并以 JavaScript Bridge 接入系统相册保存与原生分享
    \item 设计可复用样式与诗集工作流，使用户持续沉淀个人视觉语言，而非每次从零排版
\end{itemize}

\vspace{0.5em}
\noindent\textbf{学术项目}\\[-0.2em]
\rule{\textwidth}{1pt}
\textbf{自动驾驶行人感知与避障系统} \hfill 2024.01 - 2024.03\\
UIUC CS588 自动驾驶课程项目\vspace{-8pt}
\begin{itemize}[noitemsep]
    \item 构建融合目标检测与轨迹规划的自动驾驶安全系统，在检测到行人进入风险区域时控制车辆安全停车
    \item 完成立体相机与 LiDAR 的标定及多传感器融合，估计行人的三维位置、运动方向与实时轨迹
    \item 根据行人跟踪结果动态调整车辆路径与碰撞规避策略，避免静态路线无法响应移动障碍物的问题
    \item 通过真实车辆道路测试验证感知--规划--控制闭环，使系统能够随行人运动实时调整驾驶行为
\end{itemize}

\noindent\textbf{面向微视频理解的稀疏多模态序列学习} \hfill 2021.03 - 2021.05\\
本科毕业论文\vspace{-8pt}
\begin{itemize}[noitemsep]
    \item 面向微视频分类提出结合稀疏编码的多模态序列学习模型，联合建模视觉、音频与文本时序信息
    \item 使用 3 路并行 LSTM 提取各模态时序特征，并通过字典学习获得可解释的稀疏联合表示
    \item 在包含 20,000 个样本、22 个标签的数据集上与 3 个先进基线开展受控对比实验
    \item 相比基线模型，Micro-F 平均提升 45.9\%，分类准确率平均提升 53.3\%
\end{itemize}

  
% \noindent\textbf{AI驱动社交媒体生成系统
% } \hfill 2025.10\\
% 多平台支持的端到端多源数据RAG工作流\vspace{-8pt}
% \begin{itemize}[noitemsep]
%   \item 设计并实现了基于 FastAPI 的 AI驱动社交媒体内容生成系统，集成 LangChain、OpenAI
%   GPT-4o-mini 和向量数据库 Qdrant，实现 YouTube教程视频自动化转换为 Twitter 图文内容
%   \item 构建了高性能异步微服务架构，采用 PostgreSQL(SQLAlchemy 2.0 async) + Redis + Celery分布式任务队列，支持并发处理多平台内容生成请求，通过Docker Compose 实现多服务容器化编排和部署。
%  \item 开发了符合 MCP (Model Context Protocol) 标准的 YouTube 数据源服务器，实现视频转录提取、元数据分析和内容向量化，为 AI Agent 提供标准化数据接口，支持 RAG(检索增强生成) 工作流

% \item 驱动LangChain建立了多agents协作的输入 → 内容理解 → RAG增强 → 平台适配 → 输出的多源数据统一接口统一发布的工作流，为个体IP曝光、资源复用、企业营销助力
 
% \end{itemize}


% \noindent\textbf{课外竞赛}\\
% \rule{\textwidth}{1pt}
% \textbf{全国大学生数学建模竞赛} \hfill 2019.09\\
% 国家二等奖(前3\%)
% \begin{itemize}[noitemsep]
%     \item 设计分析层次过程和熵权法评估出租车决策指标
%     \item 设计蒙特卡罗模拟, 使用拟合概率曲线为出租车和乘客制定调度策略
%     \item 结合仿真和检验开发数据处理、分析和可视化程序
% \end{itemize}

% \noindent\textbf{NXP全国大学生智能汽车竞赛} \hfill 2018.01 - 2019.08\\
% 省级二等奖
% \begin{itemize}[noitemsep]
%     \item 构建能在特定区域执行巡航、避障和目标跟踪的自动驾驶汽车
%     \item 用C语言在基于ARM的微控制器上编程, 驱动发动机、摄像头和各种传感器
%     \item 应用传统滤波方法进行图像分割, 应用深度学习技术进行目标识别
%     \item 实现模糊PID控制器, 并使用独立积分和良好调整的参数对其进行优化
%     \item 实现速度高达3 m/s的平稳信标追踪, 同时灵活避免障碍物碰撞
% \end{itemize}

\end{document}
